\danceName{The Arcadian Nuptials}{longTMin}%danceName must be followed by blank line for proper command parsing

{\large Longways triple minor; $AABBCC$ adlib.\hfill \Source{Straight1784}}

\begin{music}
    \smallmusicsize%
    \parindent0mm%
    \raggedbottom%
    \stafftopmarg=1\Interligne%
    \staffbotmarg=1\Interligne%
    \nobarnumbers%
    \generalmeter{\meterfrac98}%
    \generalsignature{-1}%
    \startpiece%
%%%%%%%% PART A %%%%%%%%
    \zleftrepeat%
    \notes\qup1\Tqbu313\Tqbl535\en\xbar%
    \notes\qup1\Tqbl646\qup2\en\xbar%
    \notes\qup1\Tqbu313\Tqbl535\en\alaligne%
    \notes\Tqbl678\Tqbl987\qlp8\en\leftrightrepeat%
%%%%%%%% PART B %%%%%%%%
    \notes\qlp{10}\Tqbl{10}86\Tqbl{10}86\en\xbar%
    \notes\Tqbl789\Tqbl975\Tqbl975\en\alaligne%
    \notes\qlp{10}\Tqbl{10}86\Tqbl975\en\xbar%
    \notes\Tqbl678\Tqbl987\qlp8\en\leftrightrepeat%
%%%%%%%% PART C %%%%%%%%
    \notes\Tqbu531\Tqbu531\Tqbl358\en\alaligne%
    \notes\Tqbl642\Tqbl642\Tqbl469\en\xbar%
    \notes\Tqbu531\Tqbu531\Tqbl358\en\xbar%
    \notes\Tqbl543\Tqbu210\qup1\en\rightrepeat%
    \endpiece%
\end{music}

The 1$^{st}$ Cu: turn right hands \& ca{\s}t off one Cu: \Thaaa turn
with your left hands the Man ca{\s}t off \& the Wo: ca{\s}t up
\Thbaa foot it 3 at top \& bottom \Thaab foot it 3 contrary {\s}ides
\Thbab hands four at bottom \Thaac Right and left at top \Thbac

\HRule
\HRule
\vspace{1em}
\begin{czechDescriptionTemplate}
\begin{description}
    \item[A1]
    \item[1-2]Popis v českém jazyce píšeme na doby odpovídající hudbě, kterou máme k dispozici
    \item[3-4]Používáme ustálená spojení u typických figur, kde to dává smysl
    \item[1-4]Jinak používáme natolik přesný popis, abychom ho sami pochopili kdybychom tanec neznali
    \item[5-8]Set and turn vlevo
    \item[9-16]Ditto, set and turn vpravo
    \item[B1]
    \item[1-2] Zapisujeme tanec v pořadí, v jakém má být tančen, nepřeskakujeme části
    \item[3-4] 1.P přejde k 2.D a podá levou ruku
    \item[5-6] 1.P přejde k 3.D a podá pravou, pak levou ruku
    \item[7-8] 1.P políbí 3.D na pravou a levou tvář
    \item[9-10] 1.P a 3.D zátočka za obě o celé, 2.P a 3.P step up
    \item[11-20] dtto vše, ale provádí 1.D, na konci 2.D a 3.D step up
    \item[A2]
    \item[1-16] Sides R, set and turn, sides L, set and turn
    \item[B2] = B1
    \item[A3]
    \item[1-16] Arms R, set and turn, arms L, set and turn
    \item[B3] = B1
\end{description}
\todo{Popis tance}
\end{czechDescriptionTemplate}